
\subsection{Objectifs principaux}
\begin{itemize}
    \item \textbf{Équilibrer Performance et Énergie :} Développer un système capable de minimiser l'empreinte énergétique des serveurs tout en garantissant le respect des accords de niveau de service (SLA).
    \item \textbf{Explorer l'usage de l'IA :} Dépasser les règles de gestion fixes par l'implémentation d'agents d'Apprentissage par Renforcement (\textit{Reinforcement Learning}) capables d'apprendre des politiques de \textit{scaling} et de \textit{scheduling}.
    \item \textbf{Faire une analyse comparative :} Confronter les performances des approches des travaux précédents, des heuristiques du premier semestre (APE-H) aux modèles d'IA développés au second semestre.
\end{itemize}

\subsection{Cahier des charges technique}
\begin{itemize}
    \item \textbf{Environnement de simulation :} Utiliser et adapter un simulateur \textit{FaaS} pour reproduire une infrastructure \textit{Edge} avec placement et autoscaling de fonctions.
    \item \textbf{Création de stratégies :} Mettre en place des stratégies d'ordonnancement (\textit{scheduling}, \textit{autoscaling}, placement) afin de répondre aux contraintes de performance d'une \textit{Smart City}.
    \item \textbf{Développement IA :} Implémenter et calibrer les algorithmes DQN, PPO et une architecture Multi-Agent via l'interface standardisée \textit{Gymnasium}~\cite{brockman2016gym}.
    \item \textbf{Monitoring :} Mettre en place un suivi en temps réel des métriques (latence, énergie, \textit{cold starts}) via \textit{MatPlotLib} et \textit{TensorBoard}~\cite{abadi2016tensorflow} pour analyser les performances des modèles.
\end{itemize}
